%% paper_blind.tex
%% Blind manuscript for double-blind peer review
%% Submitted to: Systems & Control Letters
%% Author information and code repository omitted for review.

\documentclass[final,3p,times]{elsarticle}

\usepackage{amsmath,amssymb,amsthm}
\usepackage{booktabs}
\usepackage{microtype}
\usepackage[colorlinks=true,citecolor=blue,urlcolor=blue,linkcolor=black]{hyperref}

% Auto-generated by lgds.py -- do not edit by hand

\newcommand{\CondAAngleOne}{54.1}
\newcommand{\CondAAngleTwo}{24.0}
\newcommand{\CondARegretUoneTa}{0.000}
\newcommand{\CondARegretUoneTb}{0.036}
\newcommand{\CondARegretUtwoTa}{0.763}
\newcommand{\CondARegretUtwoTb}{0.000}
\newcommand{\CondARegretUjointTa}{0.026}
\newcommand{\CondARegretUjointTb}{0.026}
\newcommand{\GammaCondA}{0.026}
\newcommand{\CondBMaxAngleDeg}{0.000}
\newcommand{\CondBPertEps}{0.1}
\newcommand{\CondBPertDraws}{20}
\newcommand{\CondBPertMaxRad}{0.102}
\newcommand{\CondBPertThreshold}{0.15}
\newcommand{\HiDimFRandMean}{0.118}
\newcommand{\HiDimFPB}{0.617}
\newcommand{\HiDimFPF}{0.619}
\newcommand{\HiDimRegret}{0.0023}
\newcommand{\HiDimCoarse}{0.362}
\newcommand{\HiDimInstDep}{0.144}
\newcommand{\HiDimUpper}{1.015}


\newtheorem{theorem}{Theorem}
\newtheorem{proposition}[theorem]{Proposition}
\newtheorem{corollary}[theorem]{Corollary}
\newtheorem{definition}[theorem]{Definition}
\newtheorem{remark}[theorem]{Remark}

\newcommand{\Gr}{\mathrm{Gr}}
\newcommand{\tr}{\mathrm{tr}}
\newcommand{\R}{\mathbb{R}}
\renewcommand{\qedsymbol}{$\blacksquare$}

\journal{Systems \& Control Letters}

\begin{document}

\begin{frontmatter}

\title{Task-Dependent Optimal Linear Observers in
       Linear Gaussian Dynamical Systems}

%% Author information omitted for double-blind peer review.

\begin{abstract}
For a stationary linear Gaussian dynamical system, the rank-$r$ linear
observer minimizing squared-error prediction loss for target $Cx_{t+\tau}$
has row space equal to the top-$r$ eigenspace of the target information
matrix $M_{C,\tau} = \Sigma^{1/2}(A^\tau)^\top C^\top CA^\tau\Sigma^{1/2}$,
uniquely up to eigenspace rotation.
Two targets with distinct dominant eigenspaces admit no common optimal
rank-$r$ observer.
A common optimal observer exists if and only if the target family is
\emph{$r$-coherent}: all information matrices share a common top-$r$
eigenspace.
This condition is exact, not asymptotic or approximate.
Numerical illustrations at $n=8$, $r=2$ confirm a target-dependent case
($\Gamma=\GammaCondA>0$) and an $r$-coherent construction, both consistent
with the theorem.
\end{abstract}

\begin{keyword}
linear Gaussian dynamical systems \sep low-rank linear observers \sep
optimal observation \sep Grassmannian \sep Ky Fan theorem
\MSC[2020] 93B07 \sep 93E10 \sep 15A18
\end{keyword}

\end{frontmatter}

% -------------------------------------------------------------------------
\section{Introduction}
% -------------------------------------------------------------------------

A central question in observer design for linear dynamical systems is whether
a single rank-$r$ linear observer can be simultaneously optimal for predicting
multiple output channels $\{C_k x_{t+\tau_k}\}$, or whether the optimal
observer depends on the specific prediction target and horizon.
For finite-dimensional linear Gaussian dynamical systems (LGDS) at
stationarity, this paper gives an exact answer: the optimal rank-$r$ observer
for predicting $Cx_{t+\tau}$ is the top-$r$ eigenspace of the target
information matrix $M_{C,\tau}$, and a common optimal observer across a family
of targets exists if and only if all information matrices share a common
dominant eigenspace.

This paper is not about dynamic observer synthesis in the Luenberger--Kalman
sense \cite{Luenberger1971,OReilly1983,Kailath1980}: the observer is a fixed rank-$r$
linear map $\phi_U(x)=U\Sigma^{-1/2}x$, and performance is squared-error
prediction loss for a specified target $Cx_{t+\tau}$ at stationarity.
The static case ($A=I$, $\tau=0$) recovers standard principal-subspace
selection, but the main object here is the interaction of system dynamics,
prediction horizon $\tau$, and observer rank through the information matrix
$M_{C,\tau}$.

Reduced representations for linear systems have been studied via principal
components of Gramians \cite{Moore1981}, balanced truncation, and related
model-reduction methods \cite{Antoulas2005}.
The observability Gramian accumulates contributions from all prediction
horizons and all output directions; balanced truncation minimises
approximation error in $H_2$/$H_\infty$ system norms rather than
squared-error prediction loss for a specified $(C,\tau)$.
Sensor-selection and output-design methods target estimation or
reconstruction quality for a given objective, typically without
characterising when a fixed low-rank map is simultaneously optimal
across an entire task family.
The present note asks precisely that structural question: given distinct
prediction targets with information matrices $\{M_{C_k,\tau_k}\}$, when
is one rank-$r$ observer simultaneously optimal for all of them, and when
is target-dependent selection unavoidable?

The Grassmannian $\Gr(r,n)$ is the natural parameter space for rank-$r$
linear observers \cite{Edelman1998}, and observer quality---the trace
$\tr(U M_{C,\tau} U^\top)$---is a smooth function on this manifold.
Optimal observer selection reduces to trace maximisation on $\Gr(r,n)$,
governed by the Ky Fan variational principle \cite{Fan1949}: the maximum
is achieved by the leading eigenspace of $M_{C,\tau}$, recovering
principal-subspace selection \cite{Hotelling1933} in the static ($A=I$)
case within the classical linear state-space framework \cite{Kailath1980}.

The Ky Fan variational principle and principal-subspace optimality are
classical \cite{Fan1949,Hotelling1933}.
The single-target result (Proposition~\ref{prop:opt}) applies Ky Fan to the
target information matrix $M_{C,\tau}$, which identifies which directions of
state space bear on a given prediction task.
What is not resolved by classical eigenspace theory is the task-family
question: when does a single rank-$r$ observer remain exactly optimal across
a family of targets with possibly distinct information matrices?
The structural characterization, given in Corollary~\ref{cor:coherence}, is
$r$-coherence: its non-trivial content is the backward direction, which
shows that any commonly optimal observer must be the unique top-$r$ eigenspace
of each individual $M_k$, forcing all information matrices to share a common
dominant eigenspace.
The regret gap $\Gamma$ provides a computable certificate of unavoidable
suboptimality when $r$-coherence fails.

The result is scoped to finite-dimensional LGDS, fixed-rank linear observers,
and squared-error prediction targets.

% -------------------------------------------------------------------------
\section{System and Observer Family}
% -------------------------------------------------------------------------

\paragraph{Dynamics.}
Let
\begin{equation}
  x_{t+1} = A x_t + \xi_t, \quad \xi_t \sim \mathcal{N}(0, Q),
  \quad x_t \in \R^n,
\end{equation}
with $A$ stable (spectral radius $<1$) and $Q \succ 0$.
The unique stationary covariance $\Sigma$ satisfies
$A \Sigma A^\top + Q = \Sigma$ and is strictly positive definite
\cite{Kailath1980}; $\Sigma^{1/2}$ denotes its symmetric positive-definite
square root.

\paragraph{Observer family.}
An observer of rank $r$ ($1 \le r < n$) is parametrised by a
row-orthonormal matrix $U \in \R^{r \times n}$, $U U^\top = I_r$:
\begin{equation}
  \phi_U(x) = U \Sigma^{-1/2} x.
\end{equation}
Pre-whitening by $\Sigma^{-1/2}$ places all observers on a common isotropic
footing; in the static case ($A=I$, $\tau=0$) this reduces to principal
component analysis \cite{Hotelling1933}.
Two observers differing only by an orthogonal rotation of their row spaces are
equivalent; the observer space is the Grassmannian
$\Gr(r,n)$, a smooth compact Riemannian manifold of dimension $r(n-r)$,
with metric and geodesic structure from $O(n)$ \cite{Edelman1998}.

\paragraph{Target family.}
A prediction target is a pair $(C, \tau)$ with $C \in \R^{p \times n}$ and
integer $\tau \ge 1$:
\begin{equation}
  y_{C,\tau} = C x_{t+\tau}.
\end{equation}
Performance is total mean squared prediction error.
A target family $\mathcal{T} = \{(C_k, \tau_k)\}$ is any finite collection.

% -------------------------------------------------------------------------
\section{Theorem}
% -------------------------------------------------------------------------

\subsection{Exact Bayes risk formula}

Given observation $z = U \Sigma^{-1/2} x_t$, the MMSE predictor of
$Cx_{t+\tau}$ is linear by Gaussian conjugacy.
At stationarity, $x_{t+\tau} = A^\tau x_t + w_\tau$ where $w_\tau$
is the accumulated process noise independent of $x_t$, so
$E[x_{t+\tau}x_t^\top] = A^\tau\Sigma$.
It follows that
$\mathrm{Cov}(Cx_{t+\tau},\,z^\top)
 = CA^\tau\Sigma\cdot\Sigma^{-1/2}U^\top = CA^\tau\Sigma^{1/2}U^\top$
and $\mathrm{Cov}(z,z^\top) = U\Sigma^{-1/2}\Sigma\Sigma^{-1/2}U^\top = I_r$.
Applying the Gaussian linear MMSE formula
$\mathcal{R}(U;\,C,\tau)
 = \mathrm{tr}(\mathrm{Var}(Cx_{t+\tau}))
 - \mathrm{tr}\!\bigl(\mathrm{Cov}(Cx_{t+\tau},z^\top)
                       \mathrm{Cov}(z,z^\top)^{-1}
                       \mathrm{Cov}(z,Cx_{t+\tau})\bigr)$
gives
\begin{equation}
  \mathcal{R}(U;\,C,\tau)
    = \tr(C\Sigma C^\top) - \tr\!\bigl(U\,M_{C,\tau}\,U^\top\bigr),
\end{equation}
where the \emph{target information matrix} is
\begin{equation}
  M_{C,\tau}
    = \Sigma^{1/2}(A^\tau)^\top C^\top C A^\tau \Sigma^{1/2}
    \;\in\; \R^{n\times n}.
\end{equation}
$M_{C,\tau}$ is symmetric positive semidefinite of rank at most $\min(p,n)$.
The formula is exact.

\begin{remark}
In the static case ($A=I$, $\tau=0$),
$M_{C,0}=\Sigma^{1/2}C^\top C\Sigma^{1/2}$,
which differs from the state covariance $\Sigma$ used by PCA unless
$C=I_n$.
For $C=e_k^\top$ (predicting one component), $M_{C,0}$ is rank-$1$ with
eigenvector $e_k$ in the whitened basis: the optimal rank-$1$ observer
captures the $k$-th whitened direction regardless of the state variance in
other components, whereas PCA selects the maximum-variance direction.
In the dynamic case, the rotation by $(A^\tau)^\top$ in $M_{C,\tau}$
shifts the dominant eigenspace with~$\tau$: the same target $C$ admits a
horizon-dependent optimal observer.
\end{remark}

\subsection{Top-eigenspace optimality}

\begin{proposition}\label{prop:opt}
For fixed rank $r$, the observer minimising $\mathcal{R}(U;\,C,\tau)$ over
$\Gr(r,n)$ has row space equal to the top-$r$ eigenspace of $M_{C,\tau}$.
The minimum risk is
\begin{equation}
  \mathcal{R}^*
    = \tr(C\Sigma C^\top) - \sum_{i=1}^{r} \lambda_i(M_{C,\tau}),
\end{equation}
where $\lambda_1 \ge \cdots \ge \lambda_n$ are the eigenvalues of $M_{C,\tau}$.
Uniqueness up to rotation holds when $\lambda_r > \lambda_{r+1}$.
\end{proposition}

Degeneracy at the gap ($\lambda_r = \lambda_{r+1}$) yields a continuum of
optimal observers spanning the degenerate eigenspace; in practice, a generic
target $(C,\tau)$ produces a strict spectral gap, making degeneracy
non-generic for fixed $(A,\Sigma)$.

\begin{proof}
Minimising $\mathcal{R}$ is equivalent to maximising
$\tr(U M_{C,\tau} U^\top)$ subject to $UU^\top = I_r$.
By the Ky Fan variational principle \cite{Fan1949,HornJohnson2013}
the maximum equals
$\sum_{i=1}^r \lambda_i(M_{C,\tau})$, achieved exactly when the row space of
$U$ spans the corresponding eigenspace.
\end{proof}

The \emph{task regret} of observer $U$ for target $(C,\tau)$ is the excess
Bayes risk:
\begin{equation}\label{eq:regret}
  \mathrm{Reg}(U;\,C,\tau)
  \;:=\;
  \mathcal{R}(U;\,C,\tau) - \mathcal{R}^*(C,\tau)
  \;=\;
  \sum_{i=1}^r\lambda_i(M_{C,\tau})-\tr(U\,M_{C,\tau}\,U^\top)
  \;\ge\; 0.
\end{equation}

\begin{corollary}[Regret--angle bound]\label{cor:regretbound}
Let $M=M_{C,\tau}$ with eigenvalues $\lambda_1\ge\cdots\ge\lambda_n\ge 0$
and top-$r$ eigenspace $U^*$.
For any $U\in\Gr(r,n)$:
\begin{equation}\label{eq:regretbound}
  (\lambda_r-\lambda_{r+1})\,\|\sin\Theta\|_F^2
  \;\le\;
  \mathrm{Reg}(U;\,C,\tau)
  \;\le\;
  (\lambda_1-\lambda_n)\,\|\sin\Theta\|_F^2,
\end{equation}
where $\|\sin\Theta\|_F^2:=\|\sin\Theta(U,U^*)\|_F^2$
\textup{\cite{BjorckGolub1973}} is the squared Frobenius norm of the
principal-angle sines between $\mathrm{row}(U)$ and $\mathrm{row}(U^*)$.
\end{corollary}

\begin{proof}
Let $p_i:=\sum_{k=1}^r(v_i^\top u_k)^2$ where $v_i$ are eigenvectors of
$M$ and $u_k$ are orthonormal rows of $U$.
Then $\mathrm{Reg}(U)=\sum_{i\le r}\lambda_i(1-p_i)-\sum_{i>r}\lambda_ip_i$
and $\|\sin\Theta\|_F^2=\sum_{i\le r}(1-p_i)=\sum_{i>r}p_i$
\cite{BjorckGolub1973}.
Centering at~$\lambda_{r+1}$:
\[
  \mathrm{Reg}(U)
  =\underbrace{\textstyle\sum_{i\le r}(\lambda_i-\lambda_{r+1})(1-p_i)}_{\ge 0}
  +\underbrace{\textstyle\sum_{i>r}(\lambda_{r+1}-\lambda_i)\,p_i}_{\ge 0}.
\]
\emph{Lower bound}: drop the second sum; bound $\lambda_i-\lambda_{r+1}
\ge\lambda_r-\lambda_{r+1}$ for $i\le r$.
\emph{Upper bound}: bound the first coefficient by $\lambda_1-\lambda_{r+1}$
and the second by $\lambda_{r+1}-\lambda_n$; the two $\|\sin\Theta\|_F^2$
factors sum as written.
\end{proof}

\subsection{Target-dependent observer selection}

\begin{corollary}\label{cor:noglobal}
If targets $(C_1,\tau_1)$ and $(C_2,\tau_2)$ have information matrices
$M_1$ and $M_2$ with distinct top-$r$ eigenspaces, no single rank-$r$ observer
is simultaneously optimal for both.
\end{corollary}

\begin{proof}
Optimality for target $k$ requires the row space of $U$ to equal the top-$r$
eigenspace of $M_k$.
Distinct eigenspaces cannot be spanned by the same $U$.
\end{proof}

The corollary extends to any finite target family: if $\{M_k\}$ do not share
a common top-$r$ eigenspace, no common optimal rank-$r$ observer exists.

\subsection{\texorpdfstring{$r$}{r}-coherence: the exact exception}

\begin{definition}\label{def:rcoh}
A target family $\mathcal{T}$ is \emph{$r$-coherent} with respect to
$(A,\Sigma)$ if all matrices $\{M_{C,\tau}:(C,\tau)\in\mathcal{T}\}$
share a common top-$r$ eigenspace $V^*$.
\end{definition}

\begin{corollary}\label{cor:coherence}
A common optimal rank-$r$ observer exists if and only if $\mathcal{T}$ is
$r$-coherent.
\end{corollary}

\begin{proof}
(\emph{If})\enspace If $\mathcal{T}$ is $r$-coherent with shared top-$r$
eigenspace $V^*$, the observer with row space $V^*$ minimises
$\mathcal{R}(\cdot;\,C_k,\tau_k)$ for each $k$ by Proposition~\ref{prop:opt}.
(\emph{Only if})\enspace If $U$ is simultaneously optimal for all $k$, then
by Proposition~\ref{prop:opt}'s uniqueness (holding when
$\lambda_r(M_k)>\lambda_{r+1}(M_k)$ for all $k$), the row space of $U$
must equal the top-$r$ eigenspace of every $M_k$; hence all $M_k$ share
that eigenspace and $\mathcal{T}$ is $r$-coherent.
\end{proof}

$r$-coherence requires all targets to be functions of a common $r$-dimensional
sufficient statistic across all prediction horizons---a strong alignment
between the system dynamics and the task family.
Outside $r$-coherence, the optimal observer is target-dependent.

\begin{remark}
$r$-coherence is an algebraic equality between eigenspaces, not an
approximate or asymptotic condition.
It is the exact condition for a common optimal observer to exist.
\end{remark}

\begin{remark}
$r$-coherence is strictly weaker than simultaneous diagonalisability of
$\{M_k\}$: only the dominant $r$-dimensional eigenspace must coincide,
not the full eigenbasis.
Simultaneous diagonalisability (equivalently, pairwise commutativity
$[M_i,M_j]=0$) implies $r$-coherence for every $r$, but not conversely.
For generic $(A,\Sigma)$ and target families, the top-$r$ eigenspaces of
distinct $M_k$ are distinct, so $r$-coherence is a non-generic condition
absent structural alignment between the task family and the system.
\end{remark}

\subsection{Minimax shared observer and approximate coherence}
\label{sec:minimax}

For a task family $\mathcal{T}=\{(C_a,\tau_a)\}_{a=1}^m$, define the
\emph{minimax shared observer} and \emph{minimax regret gap}:
\begin{equation}\label{eq:mm}
  U_{\mathrm{mm}}
  = \arg\min_{U\in\Gr(r,n)}\max_a\,\mathrm{Reg}(U;\,C_a,\tau_a),
  \qquad
  \Gamma
  := \min_{U\in\Gr(r,n)}\max_a\,\mathrm{Reg}(U;\,C_a,\tau_a).
\end{equation}
$U_{\mathrm{mm}}$ is the best shared rank-$r$ observer in the worst-case sense;
$\Gamma\ge 0$ is the unavoidable worst-case cost of any single observer.

\begin{corollary}\label{cor:gammazero}
$\Gamma=0$ if and only if $\mathcal{T}$ is $r$-coherent.
\end{corollary}

\begin{proof}
(\emph{If})\enspace If $r$-coherent with shared eigenspace $V^*$:
$\mathrm{Reg}(V^*;\,C_a,\tau_a)=0$ for all $a$ by
Proposition~\ref{prop:opt}, so $\Gamma=0$.
(\emph{Only if})\enspace If $\Gamma=0$, some $U_{\mathrm{mm}}$ achieves
$\mathrm{Reg}(U_{\mathrm{mm}};\,C_a,\tau_a)=0$ for all $a$, so
$U_{\mathrm{mm}}$ is optimal for every task; the only-if direction of
Corollary~\ref{cor:coherence} then gives $r$-coherence.
\end{proof}

\begin{corollary}[Approximate coherence]\label{cor:approxcoh}
If there exists $U\in\Gr(r,n)$ with
$\|\sin\Theta(U,U_a^*)\|_F^2\le\varepsilon^2$ for all $a$, then
$\Gamma\le\Lambda\,\varepsilon^2$,
where $\Lambda:=\max_a[\lambda_1(M_a)-\lambda_n(M_a)]$.
\end{corollary}

\begin{proof}
Apply the upper bound of Corollary~\ref{cor:regretbound} to $U$ for each
task and take the maximum over~$a$.
\end{proof}

% -------------------------------------------------------------------------
\section{Numerical Illustrations}
% -------------------------------------------------------------------------

The theorem is analytically exact under its stated assumptions.
We illustrate it numerically with two finite-dimensional instances:
a target-dependent case and an $r$-coherent case.
All reported quantities are computed exactly from the analytic Bayes risk
formula and numerical optimization on the Grassmannian.

\paragraph{Setup.}
State dimension $n=8$, observer rank $r=2$.
The model is a finite-dimensional linear Gaussian dynamical system
\cite{Kalman1960}.
$A = VDV^{-1}$ with eigenvalues drawn from $U[0.3,0.85]$, random
non-orthogonal $V$ (seed 100), $Q = 0.1\,I_8$.
$\Sigma$ is computed from the discrete Lyapunov equation; $\Sigma\succ 0$
is confirmed to machine precision.
All reported quantities are computed exactly from the analytic Bayes risk
formula, eigenspace calculations, and repeated projected optimization over
$\Gr(2,8)$ (300 random initializations of projected L-BFGS-B, gradient-norm
threshold $10^{-12}$, function-value tolerance $10^{-15}$).

\paragraph{Condition A: target-dependent instance.}
{\sloppy
Targets $T_1 = ([e_1^\top;\,e_2^\top],\,1)$ and
$T_2 = ([e_5^\top;\,e_6^\top],\,5)$.\par}
With non-normal $A$, $A^1$ and $A^5$ rotate state-space directions
differently; the principal angles \cite{BjorckGolub1973} between the top-2
eigenspaces of $M_1$ and $M_2$ are $\CondAAngleOne^\circ$ and
$\CondAAngleTwo^\circ$, confirming that no single rank-2 observer is
simultaneously optimal.

The task regret \eqref{eq:regret} and minimax gap $\Gamma$
\eqref{eq:mm} are computed from the exact Bayes risk formula.
By Corollary~\ref{cor:gammazero}, $\Gamma>0$ certifies non-$r$-coherence.
The numerically computed minimax observer $U_{\mathrm{joint}}$
achieves $\Gamma=\GammaCondA>0$; balanced regrets on both targets
are a necessary condition for a minimax solution, consistent with the
global minimum having been attained.
Table~\ref{tab:condA} gives the full regret structure~\eqref{eq:regret}.

\begin{table}[tbp]
\caption{Regret matrix, Condition A ($n=8$, $r=2$, seed 100).
All values from the exact analytic Bayes risk formula.}
\label{tab:condA}
\begin{center}
\begin{tabular}{lcc}
\toprule
Observer & Regret on $T_1$ & Regret on $T_2$ \\
\midrule
$U^*_1$ (optimal for $T_1$) & \CondARegretUoneTa & \CondARegretUoneTb \\
$U^*_2$ (optimal for $T_2$) & \CondARegretUtwoTa & \CondARegretUtwoTb \\
$U_{\mathrm{joint}}$ (compromise) & \CondARegretUjointTa & \CondARegretUjointTb \\
\bottomrule
\end{tabular}
\end{center}
\end{table}

\paragraph{Condition B: structurally coherent example.}
A target family is constructed so that all $M_k$ share the same top-2
eigenspace by design: normal $A_{\mathrm{sym}} = VDV^\top$ with orthogonal
$V$, and a common $C$ whose rows are the top-2 eigenvectors of $A$, varying
only the prediction horizon $\tau \in \{1,2,3,5\}$.
Under this construction $M_{C,\tau}$ has the same dominant 2-dimensional
eigenspace for all $\tau$, so the four optimal observers are identical by the
theorem.
The maximum pairwise principal angle across all four is
$\CondBMaxAngleDeg^\circ$ to machine precision.

Under $\varepsilon = \CondBPertEps$ perturbations to $C$
(\CondBPertDraws{} random draws), the maximum principal angle across all
draws and all pairs is $\CondBPertMaxRad$~rad---well below the fragility
threshold of $\CondBPertThreshold$~rad, set at $1.5\varepsilon$.
The threshold $1.5\varepsilon$ is a heuristic calibration; the direct
numerical check confirms that the shared dominant eigenspace is robust to
this level of perturbation in $C$.
Structural coherence is perturbation-robust, distinguishing it from
accidental near-coincidences of eigenspaces that break under small
perturbations.

The numerical results are consistent with the theorem in all cases.

% -------------------------------------------------------------------------
\section{Scope}
% -------------------------------------------------------------------------

\paragraph{What is earned.}
In stationary finite-dimensional LGDS, among rank-$r$ linear observers
$\phi_U(x)=U\Sigma^{-1/2}x$ and squared-error prediction targets
$y=Cx_{t+\tau}$, no common optimal rank-$r$ observer exists unless the
target family is $r$-coherent.
$r$-coherence requires all prediction targets to be functions of a common
$r$-dimensional sufficient statistic across all horizons---a strong alignment
between the system dynamics and the task family.
Outside $r$-coherence, no common optimal observer exists, and the cost of
applying any fixed observer to a non-native target is given exactly by the
Bayes risk formula.
In observer design workflows where a single rank-$r$ map must serve multiple
prediction channels, $\Gamma$ is therefore a computable lower bound on
worst-case suboptimality under any shared observer: if $\Gamma>0$, dedicated
per-target observers are necessary to achieve optimality.
The Grassmannian $\Gr(r,n)$ is the literal domain of this optimisation:
a compact Riemannian manifold on which observer quality is a smooth function,
with $O(n)$ acting by transformation.
For fixed $C$, the computation reduces to a single $n \times n$
eigendecomposition of $M_{C,\tau}$ per horizon once $A^\tau$ is formed;
$A^\tau$ is obtained in $O(n^3 \log \tau)$ operations by repeated squaring.

\paragraph{What is not earned.}
This result concerns optimal reduced linear observation for prediction; it
does not address dynamic observer synthesis, asymptotic state reconstruction,
or feedback design.
It also does not apply to nonlinear systems, non-Gaussian noise,
nonlinear observers, or non-prediction tasks.
The breakdown is structural: for non-Gaussian noise or nonlinear dynamics,
the MMSE predictor is generally nonlinear and
Proposition~\ref{prop:opt}'s trace maximisation has no closed-form counterpart;
for nonlinear observers, the Grassmannian parametrisation and the Ky Fan
principle both cease to apply.
The causal axis---whether the observer optimal for intervention differs from
the observer optimal for prediction---is not addressed here.

\paragraph{Extensions.}
The linear Gaussian setting admits closed-form analysis because the MMSE
predictor is linear, the stationary distribution is Gaussian, and the
optimal observer reduces to a single eigendecomposition.
Whether analogous shared-eigenspace conditions characterise common optimal
observers in non-Gaussian or nonlinear systems, and whether intervention
objectives impose additional constraints on observer selection, are open
questions beyond the scope of the present result.

% -------------------------------------------------------------------------
\section*{Acknowledgements}
% -------------------------------------------------------------------------

The author received no external funding for this work.
Reproducible code is available from a publicly accessible repository
(details provided in the non-blind version of the manuscript).

\bibliographystyle{elsarticle-num}
\bibliography{refs}

\end{document}
